\section{Detector Model and Slow-Rate Boundary}\label{app:boundary}
For source $s$, the detector scores the first observed outcome to a
destination service within the contact retention horizon. For each such
outcome it updates $L\leftarrow e^{-\Delta/\tau}L+x$, where
$x=\ln(\theta_1/\theta_0)$ for success and
$x=w\ln((1-\theta_1)/(1-\theta_0))$ for failure. The default success
probabilities are $\theta_0=0.8$ and $\theta_1=0.2$; $\tau=600$\,s.
The popularity weight is $w=\max(0.25,1/(1+n))$, where $n$ counts other
remembered sources for the service. With nominal $\alpha=0.01$ and
$\beta=0.99$, the confirmation boundary is $\eta=\ln(\beta/\alpha)$,
and suspicion starts at $\eta/2$. These parameters are score settings,
not deployed error guarantees.

For periodic failed contacts with unused services, weight one, no resets
before the first alarm, and no eviction, $L_n=c+e^{-\Delta/\tau}L_{n-1}$.
Starting at zero, its limit is $c/(1-e^{-\Delta/\tau})$. With $0<c<\eta$,
the score crosses $\eta$ in finite time only when this limit exceeds $\eta$.
Fig.~\ref{fig:slow-integrated} compares the predicted boundary with nine
detector runs. Runs stop at first alarm or five hours; observed peaks on
detected runs stop at crossing. This is a deterministic implementation
check rather than a stochastic detection probability.

\begin{figure}[!t]
\centering
\includegraphics[width=\columnwidth]{\simfig{15_slow_boundary}}
\caption{Analytical and observed score near the slow-rate boundary.}
\label{fig:slow-integrated}
\end{figure}

\section{Additional Detector and Enforcement Results}
\newcommand{\simpanel}[1]{\includegraphics[width=.47\textwidth]{\simfig{#1}}}
\newcommand{\abpanel}[1]{\includegraphics[width=.47\textwidth]{\abfig{#1}}}
\begin{figure*}[!t]
\centering
\simpanel{01_F1_ci}\simpanel{02_FPR_ci}\\
\simpanel{03_scenario_TPR}\simpanel{04_scenario_FPR}\\
\simpanel{05_delay}\simpanel{06_RF_pr}
\caption{Held-out attempt simulation: operating points, scenario errors,
conditional delay, and the descriptive RF precision--recall curve.
Misses are reflected in recall, not in the delay distribution.}
\end{figure*}

\begin{figure*}[!t]
\centering
\simpanel{07_padding}\simpanel{08_sources}\\
\simpanel{09_failure}\simpanel{10_tau}
\caption{Synthetic workload and parameter sensitivity. Source counts here
are logical, not independently emulated host namespaces.}
\end{figure*}

\begin{figure*}[!t]
\centering
\simpanel{11_bound_H0}\simpanel{12_bound_H1}\\
\simpanel{13_memory}\simpanel{14_throughput}
\caption{Classical single-test reference, restart and weighted/decayed
simulations, and bounded-state diagnostics. The weighted recurrence does
not model learned popularity. Memory is traced allocation, not peak RSS;
throughput is instrumented attempt processing, not packet ingestion.}
\end{figure*}

\begin{figure*}[!t]
\centering
\abpanel{15_leaked_installation_packets}\abpanel{17_decision_to_initial_ack_s}\\
\abpanel{18_host_block_from_attack_start_s}\abpanel{19_last_block_duration_s}
\caption{Controlled enforcement diagnostics from the same 80 trials as
Fig.~\ref{fig:tradeoff}. Accepted TTL is a policy setting, not a delivery measurement.}
\end{figure*}

\section{Historical Evidence and Negative Results}\label{app:historical}
The original small-trace test found 16/16 detections and zero false alarms
over 176 benign host-runs for AIS-SPRT. TRW detected 14/16 and falsely
flagged 1/176. The historical offline CIC-IDS2017 Friday replay found
2/2 episodes detected but 2/8 benign hosts falsely flagged with the
morning-calibrated flood setting. Using the original emulation flood
setting increased the false count to 5/8. The scan and flood were observed
from one apparent source address; this is not multi-source validation.
These are historical outcomes, not held-out results of the updated
causal RF pipeline. A threshold sweep on the same CIC capture is
descriptive and does not justify a deployment threshold.

In the original four-host live loop, full-host rules appeared on all
three switches 0.68--1.39\,s after onset across four fast trials, and a
single slower trial was blocked after 20.2\,s. One of nine checks failed
because polling could not separate the path and host stages. These
measurements show that the original live loop could act, not that it met
a deadline or that a fixed-policy baseline would have acted later.

The expanded causal evaluation supersedes the older simulation comparison.
Tuned AIS-SPRT observed no confirmed false alarms but still generated
benign suspicions and missed many positives
(Tables~\ref{tab:confusion} and~\ref{tab:detectors}). RF had higher
macro F1 but nonzero false alarms, and the two methods were not matched
on false-positive rate. Keeping these negative outcomes visible is
essential when detector outputs can interrupt service.
